% =====================================================================
% BÁO CÁO ĐỒ ÁN 1 – File & Process Management
%
% Cách dùng khung:
%   - sections/report.tex là KHUNG: đã có sẵn mọi \section/\subsection.
%     Không cần sửa khung, chỉ cần thả tệp của mình vào thư mục sections/.
%   - Tệp chưa có thì báo cáo vẫn biên dịch được, chỗ đó hiện dòng chữ đỏ
%     ghi tên tệp còn thiếu. Thả tệp vào là nội dung tự xuất hiện.
%   - Tệp chỉ chứa NỘI DUNG, không viết lại tiêu đề mục (khung đã có).
%     Muốn chia nhỏ thì dùng \subsubsection.
%   - Mục 5, 7, 8 cả ba người cùng viết: mỗi người một tệp
%     <tên>-A.tex / -B.tex / -C.tex, CHỈ chứa các dòng \item ...
%     (không có \begin{itemize}); khung tự gộp thành một danh sách.
%   - Hình đặt trong images/, nên đặt tên có tiền tố theo mục,
%     ví dụ images/32-open-flow.png. Tài liệu tham khảo thêm vào refs.bib.
%
%   Tệp                          Người viết
%   01-phancong, 02-hoanthanh    B
%   31-ofile-ftable, 32-open     A
%   33-fd-compare, 34-fork       B
%   35-proc-fields, 36-scheduler C
%   41-struct, 42-files          A
%   43-procs                     C
%   44-files, 6-kiemthu          B
%   5-thietke-X,
%   7-khokhan-X, 8-chuahieu-X    mỗi người (X = A, B, C)
% =====================================================================
\documentclass[a4paper]{article}
\usepackage{vntex}
\usepackage{pdflscape}
\usepackage{rotating}
\usepackage{a4wide,amssymb,epsfig,latexsym,multicol,array,hhline,fancyhdr}
\usepackage{caption}
\usepackage{subcaption}
\usepackage{booktabs}
\usepackage{amsmath}
\usepackage{float}
\usepackage{lastpage} 
% \usepackage[lined,boxed,commentsnumbered]{algorithm2e}
\usepackage{enumerate}
\usepackage{color}
\usepackage{graphicx}					% Standard graphics package
\usepackage{array}
\usepackage{tabularx, caption}
\usepackage{multirow}
\usepackage[framemethod=tikz]{mdframed}% For highlighting paragraph backgrounds
\usepackage{multicol}
\usepackage{rotating}
\usepackage{graphics}
\usepackage{wrapfig}
\usepackage{geometry}
\usepackage{setspace}
\usepackage{epsfig}
\usepackage{tikz}
\usepackage{listings}
\usetikzlibrary{arrows,snakes,backgrounds}
\usepackage{hyperref}
\usepackage{minted}
\usepackage{indentfirst}
\usepackage{cases} 
\usepackage{listings}
\usepackage{algorithm}
\usepackage{algpseudocode}
\usepackage[english]{babel}
\renewcommand{\baselinestretch}{1.5}
\usepackage{xcolor}
\renewcommand{\figurename}{Hình}

\usepackage{tikz}
\usetikzlibrary{calc}
\usepackage{geometry}
% Header và footer được tính vào vùng nội dung (includehead/includefoot),
% nên lề 4 phía được đo từ mép ngoài của header/footer và chúng không tràn ra khỏi trang.
\geometry{a4paper, left=2cm, right=2cm, top=1cm, bottom=1.5cm,
          includehead, includefoot, headheight=55pt, headsep=10pt, footskip=30pt}

\lstset{
    language=C++,
    basicstyle=\ttfamily\footnotesize,
    numbers=left,
    numberstyle=\tiny,
    stepnumber=1,
    numbersep=5pt,
    backgroundcolor=\color{gray!10},
    keywordstyle=\color{blue}\bfseries,
    commentstyle=\color{green!50!black},
    stringstyle=\color{red},
    frame=single,
    breaklines=true,
}

%%%%%%%%
\usepackage{longtable}
\usepackage[shortlabels]{enumitem}
\usepackage{needspace}
% Cho phép giãn khoảng trắng thêm để các đoạn có \texttt dài không tràn lề phải.
\setlength{\emergencystretch}{3em}
\hypersetup{urlcolor=blue,linkcolor=black,citecolor=black,colorlinks=true} 

\pagestyle{fancy}
\fancyhead{} % clear all header fields
\fancyhead[L]{
 \begin{tabular}{rl}
    \begin{picture}(25,15)(0,0)
    \put(0,-16){\includegraphics[width=15mm, height=15mm]{images/khtn.png}}
    \end{picture}&
	\begin{tabular}{l}
		\textbf{\Large \ttfamily Trường đại học Khoa Học Tự Nhiên, ĐHQG HCM}\\
		\textbf{\Large \ttfamily Khoa Công Nghệ Thông Tin}
	\end{tabular}  	
 \end{tabular}
}
\fancyhead[R]{
	\begin{tabular}{l}
		\tiny \bf \\
		\tiny \bf 
	\end{tabular}  }
\fancyfoot{} % clear all footer fields
\fancyfoot[L]{\Large \ttfamily Hệ điều hành}
\fancyfoot[R]{\large \ttfamily Trang {\thepage}/\pageref{LastPage}}
\renewcommand{\headrulewidth}{0.3pt}
\renewcommand{\footrulewidth}{0.3pt}


%%%
\setcounter{secnumdepth}{4}
\setcounter{tocdepth}{3}
\makeatletter
\newcounter {subsubsubsection}[subsubsection]
\renewcommand\thesubsubsubsection{\thesubsubsection .\@alph\c@subsubsubsection}
\newcommand\subsubsubsection{\@startsection{subsubsubsection}{4}{\z@}%
                                     {-3.25ex\@plus -1ex \@minus -.2ex}%
                                     {1.5ex \@plus .2ex}%
                                     {\normalfont\normalsize\bfseries}}
\newcommand*\l@subsubsubsection{\@dottedtocline{3}{10.0em}{4.1em}}
\newcommand*{\subsubsubsectionmark}[1]{}
\makeatother

\definecolor{dkgreen}{rgb}{0,0.6,0}
\definecolor{gray}{rgb}{0.5,0.5,0.5}
\definecolor{mauve}{rgb}{0.58,0,0.82}

\lstset{frame=tb,
	language=Matlab,
	aboveskip=3mm,
	belowskip=3mm,
	showstringspaces=false,
	columns=flexible,
	basicstyle={\small\ttfamily},
	numbers=none,
	numberstyle=\tiny\color{gray},
	keywordstyle=\color{blue},
	commentstyle=\color{dkgreen},
	stringstyle=\color{mauve},
	breaklines=true,
	breakatwhitespace=true,
	tabsize=3,
	numbers=left,
	stepnumber=1,
	numbersep=1pt,    
	firstnumber=1,
	numberfirstline=true
}

\newcolumntype{C}[1]{>{\centering\arraybackslash}p{#1}}

% Khung chữ cho output in ra từ terminal của xv6.
\lstdefinestyle{ketqua}{
	language={},
	basicstyle=\ttfamily\scriptsize,
	backgroundcolor=\color{gray!10},
	frame=single,
	numbers=none,
	columns=fixed,
	keepspaces=true,
	showstringspaces=false,
	breaklines=false,
	aboveskip=1mm,
	belowskip=1mm
}

%%% Khung báo cáo: chèn tệp của từng thành viên %%%
\newcommand{\nguoiA}{A (Đinh Tiến Phát)}
\newcommand{\nguoiB}{B (Lê Hoàng Hiếu)}
\newcommand{\nguoiC}{C (Đặng Hồng Minh)}

% Dòng nhắc khi tệp chưa có: \chuaco{tên tệp}{người viết}
\newcommand{\chuaco}[2]{{\color{red}\textbf{[Chưa có nội dung]} Phần của #2, thêm tệp \texttt{sections/#1.tex}.}}

% \phan{tên tệp}{A|B|C}: chèn sections/<tên tệp>.tex nếu đã có.
\newcommand{\phan}[2]{%
	\IfFileExists{sections/#1.tex}%
		{\input{sections/#1.tex}}%
		{\par\noindent\chuaco{#1}{\csname nguoi#2\endcsname}\par}}

% \gop{itemize|dsso}{tên}: gộp các dòng \item trong sections/<tên>-A.tex,
% -B.tex, -C.tex thành một danh sách (dsso = danh sách đánh số).
\newlist{dsso}{enumerate}{1}
\setlist[dsso]{label=\textbf{\arabic*.}}
\newcommand{\gopmot}[2]{%
	\IfFileExists{sections/#1-#2.tex}%
		{\input{sections/#1-#2.tex}}%
		{\item \chuaco{#1-#2}{\csname nguoi#2\endcsname}}}
\newcommand{\gop}[2]{%
	\begin{#1}\gopmot{#2}{A}\gopmot{#2}{B}\gopmot{#2}{C}\end{#1}}

\begin{document}

\begin{titlepage}
\begin{center}
\textbf{\Large ĐẠI HỌC QUỐC GIA TP.HCM} \\
\textbf{\Large TRƯỜNG ĐẠI HỌC KHOA HỌC TỰ NHIÊN} \\
\textbf{\Large KHOA CÔNG NGHỆ THÔNG TIN}
\end{center}

\begin{tikzpicture}[remember picture, overlay]
    % Viền ngoài (nét đậm)
    \draw[line width=2pt] 
        ($(current page.north west) + (2cm,-2cm)$) 
        rectangle 
        ($(current page.south east) + (-1.5cm,2cm)$);
    % Viền trong (nét mảnh)
    \draw[line width=0.5pt] 
        ($(current page.north west) + (2.1cm,-2.1cm)$) 
        rectangle 
        ($(current page.south east) + (-1.6cm,2.1cm)$);
\end{tikzpicture}

\vspace{0.7cm}

\begin{figure}[h!]
\begin{center}
\includegraphics[width=6cm]{images/khtn.png}
\end{center}
\end{figure}

\vspace{-0.5cm}

\begin{center}
\begin{tabular}{c}
\multicolumn{1}{c}{\textbf{{\Large{HỆ ĐIỀU HÀNH}}}}\\
~~\\
\hline
\\
\multicolumn{1}{l}{\textbf{{\Large BÁO CÁO ĐỒ ÁN}}}\\
\\
\textbf{{\Large File \& Process Management}}\\
\\
\hline  
\end{tabular}
\end{center}

\vspace{0.5cm}

\begin{table}[h]
\begin{tabular}{rlr}
\hspace{4.5cm} Giảng viên hướng dẫn: & Phạm Tuấn Sơn\\
\hspace{4.5cm}                       & Lê Viết Long\\
\hspace{4.5cm} Sinh viên thực hiện: & Lê Hoàng Hiếu  &24120048\\
                                    & Đặng Hồng Minh  &24120090\\
                                    & Đinh Tiến Phát &24120405\\

\end{tabular}
\end{table}

\vspace{1cm}

\begin{center} {\normalsize Thành phố Hồ Chí Minh, \the\month/\the\year}
\end{center}
\end{titlepage}
%\thispagestyle{empty}


\thispagestyle{empty}

\newpage
\renewcommand{\contentsname}{Mục lục}
\renewcommand{\tablename}{Bảng}
\renewcommand{\figurename}{Hình}
\renewcommand{\refname}{Tài liệu tham khảo}

\tableofcontents

\newpage

\renewcommand{\arraystretch}{1.3} % Adjust row height
\setcounter{page}{1}



% KHUNG BÁO CÁO – không cần sửa tệp này.
% Mỗi \phan / \gop chèn tệp tương ứng trong sections/ nếu tệp đã có,
% chưa có thì hiện dòng chữ đỏ nhắc tên tệp (xem hướng dẫn đầu main.tex).

\section{Thông tin nhóm và Phân công công việc}
\phan{01-phancong}{B}

\section{Mức độ hoàn thành}
\phan{02-hoanthanh}{B}

\section{Tìm hiểu mã nguồn xv6}

\subsection{Trường \texttt{ofile[]} của \texttt{struct proc} và \texttt{file[]} của \texttt{ftable}}
\phan{31-ofile-ftable}{A}

\subsection{Trace system call \texttt{open}}
\phan{32-open}{A}

\subsection{So sánh cách dùng file descriptor trong \texttt{open()}, \texttt{dup()} và \texttt{fork()}}
\phan{33-fd-compare}{B}

\subsection{Trace system call \texttt{fork}}
\phan{34-fork}{B}

\subsection{Các trường \texttt{pid}, \texttt{state}, \texttt{context} của \texttt{struct proc}}
\phan{35-proc-fields}{C}

\subsection{Cơ chế lập lịch: \texttt{scheduler()}, \texttt{sched()}, \texttt{yield()}}
\phan{36-scheduler}{C}

\section{Cài đặt system call \texttt{sysinfo}}

\subsection{Cấu trúc \texttt{struct sysinfo} và luồng gọi}
\phan{41-struct}{A}

\subsection{Thu thập danh sách open file}
\phan{42-files}{A}

\subsection{Thu thập danh sách process và bộ đếm \texttt{nsched}}
\phan{43-procs}{C}

\subsection{Danh sách các file đã sửa đổi}
\gop{itemize}{44-files}

\section{Lựa chọn thiết kế và giả định}
\gop{itemize}{5-thietke}

\section{Kết quả kiểm thử}
\phan{6-kiemthu}{B}

\section{Khó khăn và các cách đã thử}
\gop{dsso}{7-khokhan}

\section{Những khái niệm chưa hiểu rõ}
\gop{dsso}{8-chuahieu}

\nocite{*}



\newpage    % Để phần tài liệu tham khảo qua trang mới
% \thispagestyle{empty}  % trang của phần tài liệu tham khảo để trống
\bibliographystyle{ieeetr}
\bibliography{refs}

\end{document}

